
\documentclass[12pt]{article}
\usepackage[margin=0.75in]{geometry}
\usepackage{fontspec}
\setmainfont{Latin Modern Roman}
\usepackage{microtype}
\usepackage{multicol}
\usepackage{fancyhdr}
\usepackage{titlesec}
\usepackage{enumitem}
\usepackage{xcolor}
\usepackage[hidelinks]{hyperref}
\setlength{\columnsep}{0.28in}
\setlength{\parindent}{1.1em}
\setlength{\parskip}{0.35em}
\setlength{\headheight}{15pt}
\titleformat{\section}{\large\bfseries\scshape}{}{0pt}{}
\titleformat{\subsection}{\normalsize\bfseries}{}{0pt}{}
\pagestyle{fancy}
\fancyhf{}
\fancyfoot[C]{\thepage}
\renewcommand{\headrulewidth}{0.4pt}
\newcommand{\focusbox}[1]{\par\vspace{0.4em}\noindent\begingroup\setlength{\fboxsep}{6pt}\colorbox{gray!12}{\begin{minipage}{\dimexpr\linewidth-2\fboxsep\relax}#1\end{minipage}}\endgroup\par\vspace{0.5em}}
\newcommand{\studentline}{\vspace{0.35em}\hrule\vspace{0.65em}}

\fancyhead[L]{Whately Elements of Logic}
\fancyhead[R]{Teacher Guide}
\begin{document}
\begin{center}
{\LARGE\bfseries Teacher Guide: Week 1}\par\vspace{0.25em}
{\large\scshape Logic Is Not Magic: What Logic Can and Cannot Do}\par\vspace{0.4em}
{12th Grade Long Logic Unit Based on Richard Whately's \textit{Elements of Logic}}\par\vspace{0.6em}
\rule{0.82\textwidth}{0.4pt}
\end{center}

\section*{Week 1 Overview}
This first week introduces Richard Whately's basic account of logic for a twelfth-grade long unit. Students should leave the week understanding that logic is not a machine that automatically discovers truth. Logic studies the structure of reasoning and helps expose bad inferences. The reading teaches the pattern of conclusion, stated reason, hidden assumption, and test; the argument lab then asks students to transfer that pattern to fresh examples and to the recurring Shakespeare anchors, \textit{Macbeth} and \textit{Julius Caesar}. This establishes the foundation for later weeks on terms, propositions, syllogisms, mood and figure, fallacies, ambiguity, induction, rhetoric, and verbal versus real disputes.

\section*{Essential Question}
What can logic actually do for a thinker, and what should we not expect it to do?

\section*{Student-Facing Materials}
\begin{itemize}[leftmargin=*]
\item \textit{Logic Is Not Magic: What Logic Can and Cannot Do} — student reading handout.
\item \textit{Week 1 Argument Lab: What Follows From What?} — transfer-practice handout with fresh argument tests, a repair task, an original argument task, and an exit claim.
\end{itemize}

\section*{Learning Goals}
Students will be able to:
\begin{enumerate}[leftmargin=*]
\item Define logic as the study of reasoning and argument structure.
\item Distinguish what logic can do from what it cannot do.
\item Identify conclusion, stated reason, and hidden assumption in a simple argument.
\item Explain why popularity, public honor, tradition, difficult language, confidence, or emotion do not by themselves prove a claim.
\item Repair a weak argument by narrowing the claim or strengthening the support.
\item Create and test an original argument.
\item Apply the four-step test to an initial \textit{Macbeth} or \textit{Julius Caesar} claim for the course portfolio.
\end{enumerate}

\section*{Key Vocabulary}
logic, reasoning, argument, premise, conclusion, deduction, fallacy, science, art, panacea, hidden assumption

\section*{Weekly Lesson Sequence}
\subsection*{Day 1: Opening Question and First Read}
\begin{itemize}[leftmargin=*]
\item Begin with the question: \textit{If someone is good at arguing, does that mean he is good at finding truth?}
\item Have students answer briefly in writing.
\item Introduce Whately as a nineteenth-century thinker trying to rescue logic from both overpraise and contempt.
\item Read the introduction and first half of the student reading aloud or in pairs.
\item Stop after the analogy paragraph. Ask: What does Whately think logic is being unfairly expected to do?
\end{itemize}

\subsection*{Day 2: Finish Reading and Annotate Whately's Claim}
\begin{itemize}[leftmargin=*]
\item Finish the reading.
\item Students mark three places: one sentence showing what logic can do, one showing what it cannot do, and one example that separates conclusion, stated reason, hidden assumption, and test.
\item Complete Reading Focus questions 1--6.
\end{itemize}

\subsection*{Day 3: Argument Skeleton Practice}
\begin{itemize}[leftmargin=*]
\item Introduce the pattern: conclusion, stated reason, hidden assumption, test.
\item Work through the modeled popularity/fairness example from the Argument Lab handout. Name it explicitly as the one repeated example from the reading.
\item Students complete the fresh Part II arguments individually, then compare answers with a partner.
\item Tell students that the goal is transfer: the reading already showed the method, but the lab asks whether they can use it on new arguments.
\item Emphasize that Whately's basic question is: \textit{What follows from what?}
\end{itemize}

\subsection*{Day 4: Repair and Create Arguments}
\begin{itemize}[leftmargin=*]
\item Students choose one weak argument from the lab and repair it by narrowing the claim or strengthening the support.
\item Students create an original test-case argument and identify its conclusion, stated reason, hidden assumption, and logical strength.
\item Discuss several student examples. Push students to separate proof from confidence, tradition, honor, popularity, and impressive language.
\end{itemize}

\subsection*{Day 5: Exit Claim Paragraph}
\begin{itemize}[leftmargin=*]
\item Students write the exit claim paragraph from the revised lab.
\item Require one reference to Whately's view and one example from school, literature, public speech, or everyday life.
\item Collect the student reading and argument lab as the week's work product.
\end{itemize}

\section*{Suggested Board Notes}
\begin{itemize}[leftmargin=*]
\item Logic does not discover every truth.
\item Logic studies whether a conclusion follows from premises.
\item Bad facts can still produce bad conclusions even when the reasoning looks neat.
\item A reading example may teach the method; fresh examples test whether students can transfer it.
\item Popularity, public honor, tradition, difficult language, confidence, and repetition are not the same as proof.
\item The first logical question: What follows from what?
\end{itemize}

\section*{Argument Lab Structure}
\begin{itemize}[leftmargin=*]
\item \textbf{Part I} is a model from the reading. Do not treat it as the main challenge; use it to rehearse the four-part method.
\item \textbf{Part II} contains four fresh arguments. Arguments 1--3 are deliberately weak transfer cases. Argument 4 is stronger, but only if the two students' situations are truly alike in all relevant ways.
\item \textbf{Part III} asks students to repair one weak argument. Accept either a more careful conclusion or added support that actually bears on the conclusion.
\item \textbf{Part IV} asks students to create an original argument and test their own reasoning.
\item The exit claim asks students to explain when logic helps a reader most, not merely to praise logic in general.
\end{itemize}

\section*{Reading Focus Answer Guide}
\begin{enumerate}[leftmargin=*]
\item Main claim: Logic is the science and art of reasoning; it studies argument structure and helps guard against bad deductions.
\item Can do: examine whether conclusions follow; expose skipped steps, changed terms, hidden assumptions, or bad inferences.
\item Cannot do: supply facts by itself; make foolish people wise; discover every truth; replace observation, experience, memory, imagination, or moral judgment.
\item Four steps: find the conclusion, find the stated reason, identify the hidden assumption, and test whether the conclusion really follows.
\item Answers will vary. Strong responses should use one reading example and correctly name its conclusion, reason, assumption, and test. For example, the Macbeth argument concludes that Macbeth is not responsible; its stated reason is the witches' prediction; its hidden assumption is that prediction removes responsibility; the conclusion does not necessarily follow.
\item People blamed logic because some defenders promised too much for it, so critics rejected it when it failed to perform impossible tasks.
\end{enumerate}

\section*{Argument Lab Answer Guide}
The revised lab keeps one modeled argument from the reading, then moves students into fresh test cases. Strong student answers should not merely label parts; they should explain whether the reason actually earns the conclusion.

\subsection*{Model From the Reading: Popularity and Fairness}
\begin{itemize}[leftmargin=*]
\item Conclusion: The rule must be fair.
\item Stated reason: Most students like it.
\item Hidden assumption: Whatever most students like is fair.
\item Test: The conclusion does not necessarily follow. Popularity may sometimes accompany fairness, but popularity alone does not prove fairness.
\end{itemize}

\subsection*{Argument 1: Statue and Noble Life}
\begin{itemize}[leftmargin=*]
\item Conclusion: The man must have lived nobly.
\item Stated reason: The city placed a bronze statue of him in the town square.
\item Hidden assumption: Public honor always goes to people who lived nobly, or a statue proves virtue.
\item Test: The conclusion does not necessarily follow. A statue proves that a community chose to honor someone; it does not by itself prove that the person deserved the honor.
\end{itemize}

\subsection*{Argument 2: Difficult Words and Better Reasoning}
\begin{itemize}[leftmargin=*]
\item Conclusion: The first essay must have better reasoning.
\item Stated reason: The first essay uses more difficult words than the second essay.
\item Hidden assumption: More difficult language means better reasoning.
\item Test: The conclusion does not follow. Difficult words may hide weak reasoning, and plain words may express strong reasoning.
\end{itemize}

\subsection*{Argument 3: Always Done This Way}
\begin{itemize}[leftmargin=*]
\item Conclusion: The old method must be the best method.
\item Stated reason: The class has always done the assignment this way.
\item Hidden assumption: Whatever has been done for a long time is best.
\item Test: The conclusion does not necessarily follow. Tradition may preserve wisdom, but length of use alone does not prove superiority.
\end{itemize}

\subsection*{Argument 4: Same Action, Different Punishment}
\begin{itemize}[leftmargin=*]
\item Conclusion: The rule is not just.
\item Stated reason: A just rule should treat the same action the same way, but this rule punishes one student and excuses another for the same action.
\item Hidden assumption: Relevant circumstances are truly the same; no morally important difference justifies the different treatment.
\item Test: This is a stronger argument if the cases really are alike in all relevant ways. If important circumstances differ, the argument needs more information.
\end{itemize}

\subsection*{Repair and Original Test Case}
\begin{itemize}[leftmargin=*]
\item Repaired arguments should either narrow an overclaim or add a reason that directly supports the conclusion.
\item Original student arguments should include a clear conclusion and a stated reason. Stronger responses will identify a real hidden assumption rather than repeating the stated reason.
\end{itemize}

\section*{Assessment}
Use the repaired argument, original test case, Shakespeare application, and portfolio exit paragraph as formative assessment. Strong work should:
\begin{itemize}[leftmargin=*]
\item identify conclusions and reasons accurately;
\item state hidden assumptions in complete, testable form;
\item explain whether the reason actually earns the conclusion;
\item improve a weak argument by narrowing the claim or strengthening the support;
\item avoid treating impressive language, tradition, honor, or confidence as automatic proof.
\end{itemize}

\section*{Next Week Preview}
Week 2 should move from reasoning in general to \textbf{terms}: the words that carry arguments. The next guiding question should be: \textit{Why do arguments become unstable when key words are unclear?} Keep returning to \textit{Macbeth} and \textit{Julius Caesar} so students can practice new logic tools on familiar literary ground.
\end{document}
